\section{讨论与效度}
\label{sec:discussion}
\subsection{契约带来的变化}
这一设计的核心结果在于，就绪性不能脱离具体关系来判断。一条记录只能在特定状态和策略下支持特定操作，并不携带一个与上下文无关的“有效”标志。显式保留绑定关系，使系统能够给出三类响应：采纳适用的支持证据、呈现适用的阻断项，或者指出必需的支持证据缺失。后两者都会阻止准入，但对应的后续行动不同。过期源码需要重新观察；当前失败的检查需要解决其失败；缺失的批准则需要履行已配置的批准流程。

这一设计也区分了解释与授权。历史通过记录可以解释某位工作者当时为何采取行动，即使它不能授权验收后续修订。语法边可以指导调查，而不证明语义一致性。已完成的工作项可以描述进展，却不满足验证要求。保留这些记录恰恰有用，因为它们的作用并不等同于验收。

相应的代价是增加状态与响应元数据。修订绑定、生产者键和覆盖字段会占用字节，也需要明确的失效与归约规则。预算实验展示了其中一种表现：完整源码交付可以与决策记录缺失同时出现。分层披露可按需呈现更深层的证据，但它是否改善人或智能体的决策，仍需测量。因此，评价接口是否紧凑，应同时考察精简后哪些要求不再可见，不能只看响应长度。

\subsection{效度威胁}
\paragraph{构念效度。}
Gold 源码交付与编辑输入就绪状态是某些操作的前提，不是补丁正确性的度量。嵌套谓词分析所用的参考定义本身已包含可写性，因此观察到的错误是省略该条件所产生的诊断性结果。同样，序列化决策暴露率衡量的是可供检查的信息。保留的影子评分仍针对同一组开发夹具，不能据此证明独立任务上的决策经过校准。

\paragraph{内部效度。}
固定预算顺序、debug 编译以及未受控制的主机负载，限制了对计时结果的因果解释。新的执行框架不会清空操作系统缓存。运行前后的源码检查与保留的可执行文件有助于确定所报告运行的归属，但不证明执行具有环境封闭性，也不证明每项进程依赖均已被采集。22 项反例重放使用同一保留可执行文件，但每项仅执行一次，且由开发者选定；它们并未对并发调度或潜在失败的空间进行抽样。

\paragraph{外部效度。}
语料是开发过程中编写的合成场景，并按相关的语言家族组织。必需源码异常简短，夹具规模也不能近似独立工业仓库在依赖与构建方面的多样性。重复次数不会增加独立任务数量。本文没有固定模型的问题修复研究、匹配的检索或智能体基线、独立对抗评估或操作者研究。因此，结果支持的是所测试的契约，不能证明在自主性、安全性或可用性方面普遍优于其他系统。

\paragraph{实现边界。}
形式化描述抽象了实现意图；目前没有将每条代码路径与上述公式联系起来的精化证明。文件原子性以单个文件为单位，外部编辑器可能绕过协作式锁定。修订相等所能提供的信息，也仅限于已采集的输入与指纹所覆盖的范围。这些边界明确了契约发挥作用的前提，以及哪些地方需要更强的隔离或测量机制。

\subsection{具有区分力的后续实验}
下游评估应固定模型、任务集合、预算和执行环境，并在独立仓库问题上比较相互匹配的接口。应分别记录补丁正确性、过期编辑拒绝、证据刷新开销、未解决的失败和错误验收。在此基础上，可通过机制干预，在受控故障注入下禁用区分生产者的归约或修订检查，测量假设中的验收错误是否出现。这种设计可以检验因果收益；增加当前夹具的重复运行次数无法回答这一问题。
